\documentclass[10pt,a4paper]{amsart}
\usepackage{amsmath,amssymb,amsthm,graphicx}
\usepackage{verbatim}
\newenvironment{reference}{\comment}{\endcomment}
%\newenvironment{reference}{}{}
\newenvironment{slogan}{\comment}{\endcomment}
\newenvironment{history}{\comment}{\endcomment}

% For commutative diagrams we use Xy-pic
\usepackage[all]{xy}

% We use 2cell for 2-commutative diagrams.
\xyoption{2cell}
% Theorem environments.
%
\theoremstyle{plain}
\newtheorem{theorem}[subsection]{Theorem}
\newtheorem{proposition}[subsection]{Proposition}
\newtheorem{lemma}[subsection]{Lemma}

\theoremstyle{definition}
\newtheorem{definition}[subsection]{Definition}
\newtheorem{example}[subsection]{Example}
\newtheorem{exercise}[subsection]{Exercise}
\newtheorem{situation}[subsection]{Situation}

\theoremstyle{remark}
\newtheorem{remark}[subsection]{Remark}
\newtheorem{remarks}[subsection]{Remarks}

\numberwithin{equation}{subsection}

% Macros
%
\def\lim{\mathop{\mathrm{lim}}\nolimits}
\def\colim{\mathop{\mathrm{colim}}\nolimits}
\def\Spec{\mathop{\mathrm{Spec}}}
\def\Hom{\mathop{\mathrm{Hom}}\nolimits}
\def\Ext{\mathop{\mathrm{Ext}}\nolimits}
\def\SheafHom{\mathop{\mathcal{H}\!\mathit{om}}\nolimits}
\def\SheafExt{\mathop{\mathcal{E}\!\mathit{xt}}\nolimits}
\def\Sch{\mathit{Sch}}
\def\Mor{\mathop{\mathrm{Mor}}\nolimits}
\def\Ob{\mathop{\mathrm{Ob}}\nolimits}
\def\Sh{\mathop{\mathit{Sh}}\nolimits}
\def\NL{\mathop{N\!L}\nolimits}
\def\CH{\mathop{\mathrm{CH}}\nolimits}
\def\proetale{{pro\text{-}\acute{e}tale}}
\def\etale{{\acute{e}tale}}
\def\QCoh{\mathit{QCoh}}
\def\Ker{\mathop{\mathrm{Ker}}}
\def\Im{\mathop{\mathrm{Im}}}
\def\Coker{\mathop{\mathrm{Coker}}}
\def\Coim{\mathop{\mathrm{Coim}}}

% Boxtimes
%
\DeclareMathSymbol{\boxtimes}{\mathbin}{AMSa}{"02}

%
% Macros for moduli stacks/spaces
%
\def\QCohstack{\mathcal{QC}\!\mathit{oh}}
\def\Cohstack{\mathcal{C}\!\mathit{oh}}
\def\Spacesstack{\mathcal{S}\!\mathit{paces}}
\def\Quotfunctor{\mathrm{Quot}}
\def\Hilbfunctor{\mathrm{Hilb}}
\def\Curvesstack{\mathcal{C}\!\mathit{urves}}
\def\Polarizedstack{\mathcal{P}\!\mathit{olarized}}
\def\Complexesstack{\mathcal{C}\!\mathit{omplexes}}
% \Pic is the operator that assigns to X its picard group, usage \Pic(X)
% \Picardstack_{X/B} denotes the Picard stack of X over B
% \Picardfunctor_{X/B} denotes the Picard functor of X over B
\def\Pic{\mathop{\mathrm{Pic}}\nolimits}
\def\Picardstack{\mathcal{P}\!\mathit{ic}}
\def\Picardfunctor{\mathrm{Pic}}
\def\Deformationcategory{\mathcal{D}\!\mathit{ef}}
\usepackage{cleveref}
\crefname{originallemma}{original lemma}{original lemmas}
\Crefname{originallemma}{Original lemma}{Original lemmas}
\makeatletter
\newlabel{lemma-surjection-bounded-size}{{079J}{source}{}{sourcepointer-1}{}}
\newlabel{lemma-surjection-bounded-size@cref}{{[originallemma][0][]079J}{[1][0][]source}{}{}{}}
\makeatother
\title{Bounded above acyclic subcomplexes of bounded size inside any acyclic complex}
\author{The Stacks Project Authors}
\date{stacks-project:injectives.tex:local-snapshot-20261012; GNU FDL 1.2 or later}
\begin{document}
\maketitle
\noindent\textit{Editorial note.} All delivered mathematical text is the author's own, copied from the pinned Stacks Project chapter file sites-cohomology.tex; no statement or proof step was rewritten, reordered or removed.\par\medskip
\noindent\textit{Editorial note.} The proof is the author's complete construction of both assertions, with the closing-up of the chosen surjections into a subcomplex carried out in the span.\par\medskip
\noindent\textit{Editorial note.} Every labelled reference occurring inside the counted spans is carried in assembly.external\_references as a non-load-bearing navigation pointer (same-chapter ones at their definition site, cross-chapter ones at the invocation inside the counted span), so no load-bearing mathematics is replaced by a stub. No bibliography citation occurs inside the counted spans.\par\medskip
\noindent\textit{Editorial note.} The same-name check for this card scanned every staged Stacks chapter of this repository -- all sources/stacks-*/pinned-original-source.tex files, the card's own chapter excluded -- and the labels found are recorded in dedup.same\_name\_cross\_chapter on the card.\par\medskip
\noindent\textit{Editorial note.} The preamble is the chapter's own retained minimal preamble (sources/stacks-sites-cohomology/retained-preamble.tex), a verbatim excerpt of the native preamble: the Xy-pic setup and the native theorem declarations with the mathematical macros.\par\medskip
\noindent\textit{Allowed standard prerequisites (editorial).}\begin{itemize}
\item A generator of the category. Same chapter, the existence of a generator of a Grothendieck abelian category. Supplies the object U used to define the cardinal bound.
\item The surjection of bounded size. Same chapter, Lemma lemma-surjection-bounded-size. Supplies the function c and the statement that every object receives a surjection from a coproduct of at most c many copies of the generator, which the author iterates along the complex.
\end{itemize}\medskip
\ifdefined\hypertarget\hypertarget{sourcepointer-1}{}\fi
\noindent\textit{Source reference.} Original lemma 079J: sources/stacks-injectives/pinned-original-source.tex, source lines 1787--1787. Navigation pointer invoked at proof line 1824; the counted span contains the pointer only and the cited result is not reproduced inside it, so the pointer is registered as a non-load-bearing navigation stub in the established form for Stacks material. No bibliography citation is involved.\par\medskip
\begin{lemma}
\label{lemma-acyclic-quotient-complexes-bounded-size}
Let $\mathcal{A}$ be a Grothendieck abelian category. There exists a cardinal
$\kappa$ such that given any acyclic complex $M^\bullet$ we have
\begin{enumerate}
\item if $M^\bullet$ is nonzero, there is a nonzero subcomplex
$N^\bullet$ which is bounded above, acyclic, and $|N^n| \leq \kappa$,
\item there exists a surjection of complexes
$$
\bigoplus\nolimits_{i \in I} M_i^\bullet \longrightarrow M^\bullet
$$
where $M_i^\bullet$ is bounded above, acyclic, and $|M_i^n| \leq \kappa$.
\end{enumerate}
\end{lemma}

\begin{proof}
Choose a generator $U$ of $\mathcal{A}$. Denote $c$ the function of
Lemma \ref{lemma-surjection-bounded-size}.
Set $\kappa = \sup \{c^n(|U|), n = 1, 2, 3, \ldots\}$.
Let $n \in \mathbf{Z}$ and let $\psi : U \to M^n$ be a morphism.
In order to prove (1) and (2) it suffices to prove there exists a subcomplex
$N^\bullet \subset M^\bullet$ which is bounded above, acyclic, and
$|N^m| \leq \kappa$, such that $\psi$ factors through $N^n$.
To do this set $N^n = \Im(\psi)$, $N^{n + 1} = \Im(U \to M^n \to M^{n + 1})$,
and $N^m = 0$ for $m \geq n + 2$.
Suppose for some $k \leq n$ we have constructed $N^m \subset M^m$
for all $n > m \geq k$ such that
\begin{enumerate}
\item $\text{d}(N^m) \subset N^{m + 1}$ for all $m \geq k$,
\item $\Im(N^{m - 1} \to N^m) = \Ker(N^m \to N^{m + 1})$ for
all $m \geq k + 1$, and
\item $|N^m| \leq c^{\max\{n - m, 0\}}(|U|)$ for all $m \geq k$.
\end{enumerate}
Because $M^\bullet$ is acyclic, we see that the subobject
$\text{d}^{-1}(\Ker(N^k \to N^{k + 1})) \subset M^{k - 1}$ surjects onto
$\Ker(N^k \to N^{k + 1})$. Thus we can choose $N^{k - 1} \subset M^{k - 1}$
surjecting onto $\Ker(N^k \to N^{k + 1})$ with
$|N^{k - 1}| \leq c^{n - k + 1}(|U|)$ by
Lemma \ref{lemma-surjection-bounded-size}. The proof is finished by
induction on $k$.
\end{proof}

\end{document}
